\chapter{Full Policy Tables and Oracle Decomposition}
\label{app:ceiling}

\emph{Referenced from Section~\ref{sec:ceiling}.} The main text compresses the
ceiling into two summary columns. This appendix lets a reader check what that
compression hides. First, how far each individual fixed mode sits from the
best one. Second, how the idealised chooser's advantage splits between its two
separate decisions: whether \emph{any} mode will solve the task at all, and,
among tasks that some mode does solve, which mode to pick.

\section{Per-cell policy ladder}

Table~\ref{tab:app-ladder} lists every policy evaluated on the two
VisualWebArena B0 cells; the remaining six cells follow the same structure and
are included in the released artefact
\texttt{docs/analysis/cross\_sites/router\_objective\_ordering.md}.

\begin{table}[H]
\centering\small
\caption[Full policy ladder for the two 235B cells]{Full policy ladder for the
two 235B cells: the classified-ads site
($n=224$) and the forum ($n=203$). $\Delta$ columns are against the best single
mode. \texttt{route\_only} and \texttt{triage\_only} are the two oracle halves;
\texttt{oracle\_sr\_cost} additionally breaks ties by measured cost.}
\label{tab:app-ladder}
\footnotesize\setlength{\tabcolsep}{4.5pt}
\begin{adjustbox}{max width=\textwidth}
\begin{tabular}{@{}llrrrr@{}}
\toprule
& & \multicolumn{2}{c}{\cell{Classifieds$\cdot$235B}} &
\multicolumn{2}{c}{\cell{Reddit$\cdot$235B}} \\
\cmidrule(lr){3-4}\cmidrule(lr){5-6}
Policy & kind & SR \% & cost & SR \% & cost \\
\midrule
\texttt{single:DOM}      & fixed & 17.41 & 0.06962 & 14.29 & 0.10147 \\
\texttt{single:SoM}      & fixed & 27.23 & 0.07236 & 14.78 & 0.11045 \\
\texttt{single:Vision}   & fixed & 25.00 & 0.06481 & \phantom{0}7.39 & 0.09807 \\
\texttt{single:P-text}   & fixed & 15.62 & 0.06919 & 13.30 & 0.10577 \\
\texttt{single:P-prompt} & fixed & 19.64 & 0.06853 & 12.32 & 0.10163 \\
\texttt{single:P-SoM}    & fixed & 15.62 & 0.07206 & 10.84 & 0.10814 \\
\midrule
\texttt{cheapest}        & fixed & 25.00 & 0.06481 & \phantom{0}7.39 & 0.09807 \\
\texttt{best\_sr}        & fixed & 27.23 & 0.07236 & 14.78 & 0.11045 \\
\midrule
\texttt{triage\_only}    & oracle-half & 27.23 & 0.06312 & 14.78 & 0.09998 \\
\texttt{route\_only}     & oracle-half & 43.30 & 0.06701 & 26.11 & 0.10581 \\
\texttt{oracle\_sr}      & oracle & 43.30 & 0.06259 & 26.11 & 0.09794 \\
\texttt{oracle\_sr\_cost}& oracle & 43.30 & 0.05777 & 26.11 & 0.09534 \\
\midrule
\texttt{cascade\_cost\_first} & oracle (perfect det.) & 43.30 & 0.31896 & 26.11 & 0.54588 \\
\bottomrule
\end{tabular}
\end{adjustbox}
\end{table}

\FloatBarrier
\section{The cost tie-break, isolated}

\texttt{oracle\_sr} and \texttt{oracle\_sr\_cost} have identical success and
identical treatment of unsolvable tasks, so the entire gap between them is the
tie-break rule applied to tasks with more than one solver. On
\cell{cls$\cdot$B0} that is 68 of 224 tasks, and switching the tie-break from
list order to measured cost saves \$0.01589 on each, reducing mean cost from
\$0.06259 to \$0.05777 ($-7.7\%$). The corresponding figures are 36 of 203 tasks
and $-2.7\%$ on \cell{red$\cdot$B0}.

This is worth isolating because it quantifies a defect discussed in
Section~\ref{sec:mechanism}: the mode ordering used to break ties is documented
as ascending in cost and is not, so a labelling pipeline that trusts it is
systematically choosing a more expensive successful mode than necessary.

\FloatBarrier
\section{Cascade attempt histograms}

The cost-first cascade of Section~\ref{sec:overhead} escalates through modes in
ascending measured cost until a success is detected. Its attempt histograms show
where the cost goes:

\begin{itemize}
\item \cell{cls$\cdot$B0}: $\{1{:}56,\ 2{:}25,\ 3{:}3,\ 4{:}5,\ 5{:}2,\ 6{:}133\}$
\item \cell{red$\cdot$B0}: $\{1{:}15,\ 2{:}23,\ 3{:}5,\ 4{:}4,\ 5{:}2,\ 6{:}154\}$
\item \cell{wa\_red$\cdot$B0}: $\{1{:}28,\ 2{:}9,\ 3{:}10,\ 4{:}1,\ 5{:}4,\ 6{:}52\}$
\end{itemize}

The 6-attempt bucket is tasks no mode solves: the cascade pays for all six
before learning this. It is 59.4\%, 75.9\% and 50.0\% of tasks respectively,
which is the entire explanation of the $+316$ to $+463\%$ cost.
